% Generated draft: requires mathematical and source review.
\documentclass[11pt]{article}
\usepackage[margin=25mm]{geometry}
\usepackage{fontspec}
\setmainfont{Latin Modern Roman}
\usepackage{amsmath,amssymb,mathtools,unicode-math}
\setmathfont{Latin Modern Math}
\usepackage{enumitem,needspace,xurl,hyperref}
\setlength{\parindent}{0pt}
\setlength{\parskip}{5pt}
\setlength{\emergencystretch}{3em}
\newcommand{\N}{\mathbb{N}}
\newcommand{\Z}{\mathbb{Z}}
\newcommand{\Q}{\mathbb{Q}}
\newcommand{\R}{\mathbb{R}}
\newcommand{\C}{\mathbb{C}}
\newcommand{\F}{\mathbb{F}}
\newcommand{\A}{\mathbb{A}}
\newcommand{\PP}{\mathbb{P}}
\newcommand{\E}{\mathbb{E}}
\newcommand{\eps}{\varepsilon}
\newcommand{\dd}{\,\mathrm{d}}
\begin{document}
\section*{UnsolvedMath ID: 6700019}
\textbf{Source problem code:} catalogue entry
\subsection{Definitions and mathematical statement}
A Riemannian manifold is flat if its Riemann curvature tensor vanishes identically. The extracted item asks, tentatively, to determine the assumptions on \(\pi_1(X)\), depending on a flat manifold \(X_{\mathrm{fl}}\), that are needed for ``this conclusion'', for all flat \(X_{\mathrm{fl}}\). The conclusion and geometric hypotheses are absent from the item itself.

The adjacent local source record, permanent ID 6700018, gives the following context: \(X_{\mathrm{fl}}=\R^n/\Gamma\) is complete and flat, initially with \(\Gamma\) acting by translations; \(X\) is complete, has \(\operatorname{Scal}_X\geq0\), and is isometric to \(X_{\mathrm{fl}}\) at infinity; the question is whether \(X\) must be flat. Interpreting isometry at infinity as an isometry
\[
 X\setminus K\ \cong\ X_{\mathrm{fl}}\setminus K_{\mathrm{fl}}
\]
for suitable compact subsets, the intended extension appears to be: for each complete flat model \(X_{\mathrm{fl}}\), including those with nontrivial linear holonomy, determine suitable fundamental-group hypotheses under which every such \(X\) is flat.

This is a contextual reconstruction, not a complete theorem supplied by the present item. It is not stated whether the desired assumptions involve only the abstract group \(\pi_1(X)\), its relation to \(\pi_1(X_{\mathrm{fl}})\), or homomorphisms induced by the common region at infinity. Nor does ``needed'' specify necessity and sufficiency. The proposed conclusion is vanishing curvature of \(X\), not necessarily an isometry \(X\cong X_{\mathrm{fl}}\).
\subsection{Short English statement}
Determine which fundamental group conditions force a complete manifold with nonnegative scalar curvature and a prescribed flat geometry at infinity to be flat, for arbitrary flat models. The item refers to missing context; the adjacent local record supplies a provisional interpretation.
\subsection{Sources}
\begin{itemize}
\item Local UnsolvedMath registry, permanent ID 6700019, source code catalogue entry; Misha Gromov, 2017; update 2026-08-06. Adjacent local record ID 6700018, source code AMR-066-0018, was consulted solely to identify the missing antecedent.
\item Gromov, \emph{101 Questions, Problems and Conjectures around Scalar Curvature} (2017), printed tag [?19], occurrence 19, PDF page 24; the contextual item [?18] is also on page 24: \url{https://www.ihes.fr/~gromov/wp-content/uploads/2018/08/101-problemsOct1-2017.pdf}. Registry access date 2026-07-29; normalized PDF text extraction; public author-hosted incomplete PDF; present status and reuse terms require review.
\item The research summary mentions essentiality, aspherical cases, polynomial growth, and a Burago--Ivanov flat torus theorem without an exact citation or a theorem matching this question. These reported connections do not supply the missing hypotheses.
\end{itemize}
\end{document}
